\scanpage{305}
\mspage{46}
mais peut-être plus grand. (P. ex. si $\tau=\tau'$, $f=\mathrm{id}$,
\deleted{\unsure{cela contient}} $E'=I_{\tau'}$, qui est donc str. plus grand que $I_{\tau'\,\unsure{\mathrm{r\acute ed}|\mathrm{red}}}$
si $\tau'$ n'est pas réduit). Pour \deleted{tout} $i\in E'$,
considérons $f^{-1}(\{i\})\subset I_\tau$. \inserted{Comme $f$ croissant,} c'est un intervalle
de l'ens \deleted{\unsure{tot.}} ordonné $I_\tau$, donc c'est un sous-intervalle
de \unsure{la|l'} ``\unsure{catégorie|catigo|catég}'' intervalle $I_\tau$. Comme
tel, il a une signature $\tau_i\in\mathrm{Ty}$. On a donc
une application
\[
 \unsure{\eta|\upsilon}_f\deleted{\unsure{=}}\ \text{ou }\eta\colon E'\longrightarrow\mathrm{Ty}.
\]
Je dis que $(\tau,f)\in\operatorname{Ob}\,\underline{\mathrm{Ty}}_{/\tau'}$ est déterminé de
façon unique par ce couple
\[
 (E',\eta)\ \left|\begin{aligned}
 E'&\in\unsure{\mathscr P|\mathcal P}(I_{\tau'})\unsure{=|\colon}\mathscr P([0,n'])\\
 \eta&\colon E'\longrightarrow\mathrm{Ty}
 \end{aligned}\right.\quad,
\]
De plus, celui-ci est soumis à
\deleted{aucune} \inserted{la seule restriction} que \deleted{$E'$}
\[
 E'\supset\mathrm{Cr}(\tau')=\substack{\text{ens. des sommets critiques}\\\text{de }I_{\tau'}\, .}
\]
(ou ce qui revient au $=$, que $E'$ contient $0_{\tau'},1_{\tau'}$
et que deux éléments consécutifs de $E'$ sont
reliés par une flèche de $I_{\tau'}$)\unsure{;|:} \deleted{\unsure{le}}
\[
 \left\{\begin{aligned}
 \operatorname{Ob}\big(\underline{\mathrm{Ty}}_{/\tau'}\big)&\simeq\text{ens des couples }(E',\eta),\ \text{avec}\\
 &E'\subset[0,\underset{\substack{\parallel\\\operatorname{long}\tau'}}{n'}],\quad E'\supset\mathrm{Cr}(\tau')\\
 &\eta\colon E'\longrightarrow\mathrm{Ty}
 \end{aligned}\right.
\]
Soit
\[
 0=t_0<t_1<\ldots<t_N=n'
\]
la suite des éléments de $E'$, soit pour \tikz[baseline=(q.base)]\node[draw,ellipse,inner sep=2pt](q){$i\in[1,N]$};
\[
 \left\Vert\begin{aligned}
 \unsure{l_i|h_i}&=t_i-t_{i-1} &\unsure{\neq1|\geq1}\\
 \varepsilon_i&\quad\text{signature de }[t_{i-1},t_i]\ \text{dans }I_{\tau'}\\
 &\deleted{\unsure{I_{\tau}\ \eta(\ }} 
 \end{aligned}\right.
\]
% archive sheet number: 304
